\chapter{Kesimpulan dan Saran}

\section{Kesimpulan}

Penelitian ini menganalisis dan membandingkan dampak integrasi PLTB berbasis
DFIG-WT dan PMSG-WT terhadap \textit{small signal stability} (SSS) sistem tenaga
listrik pada sistem uji IEEE \textit{nine-bus} menggunakan analisis modal
berbasis \textit{eigenvalue} di DIgSILENT PowerFactory 2024.

\begin{enumerate}

\item Integrasi DFIG-WT menurunkan \textit{damping ratio} dan margin
kestabilan beban. Penggantian generator sinkron konvensional oleh DFIG-WT
menurunkan DR mode dominan dari 7,22\% (kondisi awal) menjadi sekitar 4,04\%,
baik ketika DFIG-WT menggantikan G3 (Skenario 2.1, 85 MW) maupun G2 (Skenario
2.2, 163 MW). Lokasi dan kapasitas penggantian tidak memengaruhi DR mode dominan
(4,04\% vs 4,05\%). Analisis \textit{participation factor} pada skenario
peningkatan beban menunjukkan bahwa variabel $\varphi_m$ (sudut rotor internal
DFIG-WT) mendominasi mode tidak stabil, mengonfirmasi bahwa dinamika internal
DFIG-WT memicu ketidakstabilan. Integrasi DFIG-WT juga mempersempit margin kestabilan beban secara signifikan:
batas ketidakstabilan Beban A turun dari 247\% menjadi 83\% (Skenario 4.1), dan
Beban B dari 373\% menjadi 123\% (Skenario 4.2). Ketidakstabilan yang terjadi
bersifat aperiodik dengan DR $-$100\% pada kedua \textit{eigenvalue} riil positif,
berbeda dari ketidakstabilan osilatoris pada skenario konvensional (DR $-$1,17\%
pada Skenario 1.2 dan $-$2,45\% pada Skenario 1.3), menunjukkan mekanisme
kehilangan kestabilan yang lebih cepat dan tanpa osilasi peringatan.

\item Integrasi PMSG-WT mempertahankan atau meningkatkan \textit{damping
ratio}, dengan margin kestabilan beban yang bergantung pada posisi. Penggantian
generator sinkron oleh PMSG-WT tidak menurunkan DR: Skenario 3.1 (menggantikan
G3) menghasilkan DR 7,25\% dan Skenario 3.2 (menggantikan G2) menghasilkan DR
8,73\%, keduanya lebih tinggi dari kondisi awal 7,22\%. Lokasi penggantian
generator oleh PMSG-WT memengaruhi DR, di mana penggantian G2 (kapasitas lebih
besar, 163 MW) justru menghasilkan redaman lebih tinggi. Analisis
\textit{participation factor} mengonfirmasi bahwa variabel dari model PMSG-WT
tidak muncul sebagai kontributor dominan pada mode osilasi, karena
\textit{fully-rated converter} memisahkan dinamika mesin dari jaringan. Margin kestabilan beban dengan PMSG-WT secara konsisten lebih baik dibandingkan
DFIG-WT: batas ketidakstabilan Beban A pada Skenario 5.1 (PMSG-WT menggantikan
G3) adalah 165\%, lebih dari dua kali lipat Skenario 4.1 (DFIG-WT menggantikan
G3, 83\%); batas ketidakstabilan Beban B pada Skenario 5.2 (PMSG-WT menggantikan
G2) adalah 284\%, juga lebih dari dua kali lipat Skenario 4.2 (DFIG-WT
menggantikan G2, 123\%). Nilai DR pada batas ketidakstabilan adalah $-$0,60\%
(Skenario 5.1, osilasi sangat lambat dengan $f_d = 0{,}006$ Hz) dan $-$0,075\%
(Skenario 5.2, osilasi \textit{inter-area} dengan $f_d = 0{,}340$ Hz).

\item DFIG-WT dan PMSG-WT memberikan dampak yang berlawanan terhadap
SSS, dengan perbedaan mendasar pada \textit{eigenvalue}, DR, dan
\textit{participation factor}. DFIG-WT menurunkan DR mode dominan dan
mempersempit margin kestabilan beban di semua skenario, sedangkan PMSG-WT
mempertahankan atau meningkatkan DR dengan margin kestabilan yang bergantung pada
posisi penggantian terhadap beban. Seluruh mode dominan dalam kondisi stabil
tergolong \textit{local area mode} (1,38--1,86 Hz), tanpa karakteristik
\textit{inter-area}. Jumlah \textit{eigenvalue} meningkat seiring penambahan
PLTB: 16 mode (kondisi awal), 20 mode (PMSG-WT), dan 30 mode (DFIG-WT),
mencerminkan kompleksitas model kontrol DFIG-WT yang lebih tinggi. Ketika sistem
mencapai ketidakstabilan, variabel dominan berbeda: $\varphi_m$ DFIG-WT
mendominasi pada Skenario 4.1 dan 4.2, sedangkan $\omega$ G$_1$ mendominasi pada skenario PMSG-WT (5.1) maupun skenario
konvensional (1.2, 1.3), dan $x_{s2}$ PLL PMSG-WT mendominasi pada Skenario 5.2,
mengonfirmasi bahwa \textit{fully-rated converter} PMSG-WT mencegah dinamika mesin
itu sendiri berpartisipasi langsung dalam mode tidak stabil.

\end{enumerate}

\section{Saran}

Empat aspek berikut dapat dikembangkan pada penelitian selanjutnya berdasarkan
batasan penelitian dan temuan di Bab~IV.

\begin{enumerate}

\item Penambahan AVR dan PSS pada generator sinkron konvensional.
Penelitian ini tidak menggunakan AVR maupun PSS pada generator konvensional yang
tersisa. Penambahan kedua komponen tersebut, dengan berbagai kombinasi tipe
(misalnya EXDC2, EXAC4, STAB1, PSS2A), berpotensi meningkatkan DR mode dominan
yang pada skenario DFIG-WT sudah mendekati ambang 5\%. Penelitian lanjutan dapat
mengevaluasi kombinasi AVR dan PSS terbaik untuk mempertahankan SSS sistem saat
PLTB terhubung.

\item Penerapan kontroler peredam osilasi pada DFIG-WT. Analisis
\textit{participation factor} menunjukkan bahwa variabel $\varphi_m$ DFIG-WT
mendominasi mode tidak stabil pada kondisi peningkatan beban, mengindikasikan
bahwa dinamika internal DFIG-WT menjadi sumber ketidakstabilan. Penerapan
\textit{supplementary damping controller} atau \textit{wind power stabilizer}
pada kontroler DFIG-WT merupakan langkah lanjutan untuk meredam mode tersebut.

\item Variasi tingkat penetrasi dan kecepatan angin. Penelitian ini
hanya mengkaji satu tingkat penetrasi per skenario dengan asumsi kecepatan angin
konstan pada kondisi nominal. Variasi penetrasi parsial (misalnya 25\%, 50\%,
75\% kapasitas nominal) serta kondisi angin rendah dan angin tinggi akan memetakan
batas operasi aman DFIG-WT dan PMSG-WT terhadap SSS.

\item Validasi pada sistem tenaga praktis dan perluasan jenis kajian
kestabilan. Sistem uji IEEE \textit{nine-bus} adalah sistem akademis berskala
kecil yang tidak mewakili kompleksitas sistem interkoneksi nyata. Penerapan
metodologi yang sama pada sistem kelistrikan nyata, seperti sistem interkoneksi
Sulawesi atau sistem Jawa-Bali sesuai RUPTL PLN 2025--2034, akan menguji apakah
kesimpulan ini berlaku pada skala interkoneksi penuh. Analisis dalam penelitian
ini terbatas pada SSS; penambahan analisis kestabilan transien akan mengungkap
respons sistem terhadap gangguan besar saat PLTB terhubung.

\end{enumerate}


